% GENERATED FILE. Edit canonical lessons, then run npm run generate:books.
% canonical-source-hash: fnv1a64:98ae575fd9c0a1be
% canonical-lessons: IT-C13-nero-bianco, IT-C13-rosso-blu

\chapter{Colours}
\label{ch:it-colours}
\hlchaptermodality{13}

\begin{chapteropening}
\textbf{By the end of this chapter:} \emph{I can name black, white, red and blue in Italian and say which of them Italian borrowed rather than inherited.}

Say nero, bianco, rosso and blu, add azzurro for gli Azzurri, and sort the four: nero and rosso inherited, bianco and blu Germanic loans, azzurro out of Arabic lāzaward.
\end{chapteropening}

\section[nero, bianco]{nero, bianco — one Roman, one Lombard}

\label{lesson:IT-C13-nero-bianco}



\begin{quote}
Italy's two most basic colours came from \textbf{two different peoples}. One is Rome's own word. The other arrived with the Germanic tribes who settled the peninsula after Rome fell.
\end{quote}

\begin{cousinweb}
\textbf{nero} (\textquotedblleft{}black\textquotedblright{}) $\leftarrow$ Latin \emph{\textbf{niger}}. Straight inheritance, barely changed: \emph{niger} $\to$ \emph{nigrum} $\to$ \emph{nero}. The same \emph{niger} gives English \emph{denigrate}, \textquotedblleft{}to blacken.\textquotedblright{}
\end{cousinweb}

\begin{cousinweb}
Latin's white was \emph{\textbf{albus}}. Italian doesn't use it. Instead: \textbf{bianco}, from Germanic \emph{\textbf{blank}} — \textquotedblleft{}\textbf{shining, gleaming}\textquotedblright{} — most likely carried in by the \textbf{Lombards}, the Germanic people who ruled northern Italy for two centuries and left their name on \textbf{Lombardia}.

The same word became French \emph{blanc} and Portuguese \emph{branco}. Germanic lending \textbf{into} Romance is the reverse of the usual direction, and here the loan won so completely that Latin's own word was pushed out of the colour slot altogether.
\end{cousinweb}

\begin{cousinweb}
\noindent
\begin{tabularx}{\linewidth}{@{}>{\raggedright\arraybackslash}X>{\raggedright\arraybackslash}X>{\raggedright\arraybackslash}X@{}}
\toprule
\textbf{word} & \textbf{meaning} & \textbf{the white inside} \\
\midrule
\textbf{alba} & dawn & the sky turning white \\
\textbf{albume} & egg white & the white part \\
\textbf{Albania}, \emph{Alpi} & (place names) & \textquotedblleft{}the white ones\textquotedblright{} \\
\emph{album} (Eng.) & album & a blank white tablet \\
\bottomrule
\end{tabularx}

Latin's white is still here — just doing other jobs.
\end{cousinweb}

\subsection*{Your turn}
Say these aloud:

\begin{itemize}
\raggedright
  \item \textquotedblleft{}nero, bianco\textquotedblright{}
  \item \textquotedblleft{}il vino bianco, il vino rosso\textquotedblright{} — Ch.11's wine, coloured
  \item \textquotedblleft{}il pane bianco\textquotedblright{} — Ch.11's bread
  \item \textquotedblleft{}l'\textbf{alba}\textquotedblright{} — and hear Latin's forgotten white
\end{itemize}

\begin{tcolorbox}[breakable,colback=teal!4,colframe=teal!35!black,title={Before you move on}]
Which colour is Rome's own? (\textbf{Nero} $\leftarrow$ \emph{niger}.) Where does \textbf{bianco} come from? (\textbf{Germanic \emph{blank}}, \textquotedblleft{}shining\textquotedblright{} — likely via the \textbf{Lombards}.) Why is that surprising? (Latin normally lends \textbf{to} Germanic, not the reverse.) Where did \emph{albus} survive? (In \textbf{alba}, \textquotedblleft{}dawn,\textquotedblright{} and \textbf{albume} — no longer a colour.) Next: \textbf{rosso} and \textbf{blu}, and one more traveller.
\end{tcolorbox}
\section[rosso, blu]{rosso, blu — and the blue Italy plays in}

\label{lesson:IT-C13-rosso-blu}



\begin{quote}
One very old word, one more Germanic loan — and a third blue that came all the way from Persia by way of Arabic.
\end{quote}

\begin{cousinweb}
\textbf{rosso} $\leftarrow$ Latin \emph{\textbf{russus}} (\textquotedblleft{}red\textquotedblright{}), from PIE \emph{\textbf{h\textsubscript{1}rewd\textsuperscript{h}-}}. That root is one of the oldest colour words we can reconstruct:

\noindent
\begin{tabularx}{\linewidth}{@{}>{\raggedright\arraybackslash}X>{\raggedright\arraybackslash}X@{}}
\toprule
\textbf{language} & \textbf{word} \\
\midrule
Italian & \textbf{rosso} \\
French & \textbf{rouge} ($\leftarrow$ \emph{rubeus}) \\
English & \textbf{red}, \emph{rust}, \emph{ruby} \\
German & \textbf{rot} \\
\bottomrule
\end{tabularx}

\emph{Rosso} and \emph{red} are \textbf{cousins}, not borrowings — separated for millennia. Note the \textbf{double s}: hold it. \emph{Roso} and \emph{rosso} are different words.
\end{cousinweb}

\begin{cousinweb}
\textbf{blu} $\leftarrow$ Germanic \emph{\textbf{blāo}}, exactly like French \emph{bleu}. Chapter pattern confirmed: \textbf{bianco} and \textbf{blu}, Italian's white and blue, are both \textbf{Germanic imports}.
\end{cousinweb}

\begin{cousinweb}
Italian has a second blue, and it's the famous one: \textbf{azzurro}, from Arabic \emph{\textbf{lāzaward}} (\textquotedblleft{}lapis lazuli\textquotedblright{}), itself from Persian \emph{\textbf{lāžward}}. The initial \emph{l-} was swallowed — heard as an article and dropped — leaving \emph{azur/azzurro}. The same journey gave Spanish and Portuguese \textbf{azul}, French \textbf{azur}, and English \textbf{azure}.

Italy's national teams are \textbf{gli Azzurri}, \textquotedblleft{}the blue ones\textquotedblright{} — named, at the end of a very long chain, after a \textbf{blue stone mined in Afghanistan}.
\end{cousinweb}

\begin{grammarlens}[title={What you've built: the colour paths}]
\noindent
\begin{tabularx}{\linewidth}{@{}>{\raggedright\arraybackslash}X>{\raggedright\arraybackslash}X@{}}
\toprule
\textbf{Italian} & \textbf{source} \\
\midrule
nero & Latin \emph{niger} \\
\textbf{bianco} & \textbf{Germanic} \emph{blank} \\
rosso & Latin \emph{russus} (ancient PIE root) \\
\textbf{blu} & \textbf{Germanic} \emph{blāo} \\
\textbf{azzurro} & \textbf{Arabic} \emph{lāzaward} \\
\bottomrule
\end{tabularx}
\end{grammarlens}

\subsection*{Your turn}
Say these aloud:

\begin{itemize}
\raggedright
  \item \textquotedblleft{}rosso, blu, azzurro\textquotedblright{}
  \item \textquotedblleft{}il vino rosso\textquotedblright{} — Chapter 11's wine
  \item the cousins — \textquotedblleft{}rosso · red · rot · rouge\textquotedblright{}
  \item \textquotedblleft{}gli Azzurri\textquotedblright{} — and think of the stone
\end{itemize}

\begin{tcolorbox}[breakable,colback=teal!4,colframe=teal!35!black,title={Before you move on}]
Is \emph{rosso} related to English \emph{red} by borrowing or descent? (\textbf{Descent} — PIE \emph{\textbf{h\textsubscript{1}rewd\textsuperscript{h}-}}.) Which two Italian colours are Germanic loans? (\textbf{Bianco} and \textbf{blu}.) Where does \textbf{azzurro} come from? (\textbf{Arabic \emph{lāzaward}}, \textquotedblleft{}lapis lazuli\textquotedblright{} — via Persian; the \emph{l-} was lost as if it were an article.) What are Italy's teams called? (\textbf{Gli Azzurri}.) Next chapter: attaching these colours to the nouns you know.
\end{tcolorbox}
